\chapter{Code Reuse}
\label{sec:func}

% motivation: in the last chapter we covered branching and loops for flow control, this gave us mechanisms for selective and repeated execution of statements, the repetition constructs allows us to repeatedly call some lines of code from the same location in code, that is every time we get to a specific part of a loop we execute these lines, but what if there are two places where we want to execute those lines? We certainly could duplicate those lines, but that solution doesn't scale, you will soon see situation where there are tens of such locations and accept that there can easily be tens of thousands, technically we could have that many copies but that would be expensive in terms of memory and -- more importantly -- impose a severe maintenance burden on us, what happens when we want to extend that functionality or find a bug? Then we need to find all of these potentially thousands of locations and update them, we need tooling for that
In the last chapter, we covered branching and loops as tools for control flow. This gave us mechanisms for selective and repeated execution of statements. The repetition constructs allows us to repeatedly execute some segment of our code from the same location. That is, every time we get to a specific part of a look, we execute these lines. But what if there were two places where we wanted to execute those very same lines? We certainly could duplicate those lines and call it a day, but that solution doesn't \idx{scale}{Scalability}. You will soon experience situations where you yourself have tens of such locations, and accept that others will have situations where they have to deal with tens of thousands of such locations (some of which are outside of their control). Technically, we could have that many copies, but that would be expensive in terms of memory and -- more importantly -- impose a severe \idx{maintenance}{Maintainability} burden on us: What happens when we want to extend that functionality or find a \idx{bug}{Bug} in it. Then we need to find all of these potentially thousands of locations and update them individually. This is a task that humans are ill-equipped to deal with. For us to have any certainty of all of these updates being performed without introducing new bugs we would need the support of \idx{tooling}{Tooling}.

% solution: that tooling turns out not to be unnecessary, in fact seasoned programmers try to avoid the code duplication problem most of the time, instead they package commonly used functionality in something called a function, a function is a parameterized piece of code that typically sits somewhere outside your programs main flow, it has a well-defined interface and a descriptive name, execution can be temporarily transferred to it from any location in your code, we refer to this a 'calling the function'
That tooling turns out not to be necessary. In fact, seasoned programmers usually try to avoid the \idx{code duplication}{Duplicate code} problem. Instead, they package commonly used functionality in something called a \idx{function}{Function}. A function is a parameterized piece of code that typically sits somewhere outside your programs main flow. If has a well-defined \idx{interface}{Interface} and a descriptive \idx{name}{Name}. Execution can be temporarily transferred to it from any location in your code. We refer to this as \textsl{calling the function}. This is essentially what \idx{procedural programming}{Procedural programming} is.

% code structure
Functions -- like many of the concepts that we will introduce later on -- give \idx{structure}{Code structure} to our code in ways that allows a human to navigate it. Instead of having one large program of thousands or even millions of lines of code that we need to have an overview of all the time, we get to a situation where we can often choose to focus on an individual function of perhaps 10-50 lines. With our limited capacity for handling large amounts of information, such structure is necessary for us to \idx{scale}{Scalability!Complexity} the increase of complexity.

\section{Mathematical Functions}

% definition
In calculus, a function represents a mapping from the set of real numbers to the set of real numbers. That mapping is usually expressed through a mathematical formula and can be expressed as on a 2d coordinate system. Some functions operate in a higher dimensionality. For instance, some function may map two of the sets of real numbers to the set of real numbers. Such a function could be illustrated in 3d.

% redefinition: by taking a few liberties we can define a mathematical function as a mapping from a list of values to a single value, the list of input values are matched against and assigned to the symbolic names of the function definition, the mapping is expressed -- if we put on our programming glasses -- as an expression over these symbolic names
By taking a few liberties, we can redefine a concrete mathematical function as a mapping from a list of values to a single value. The list of input values are matched against and assigned to the symbolic names of the function definition. The mapping itself is expressed -- if we put on our programming glasses -- as an expression over these symbolic names.

\subsection{Linear Functions}
\label{sec:func:linear}

In calculus, a linear function that, when illustrated forms a non-vertical straight line in a 2d coordinate system and follows the form of:
\begin{eqnarray*}
  f(\textcolor{purple}{x}) = a \cdot \textcolor{purple}{x} + b
\end{eqnarray*}

% description
It is a named formula that is parameterized over $x$. Concrete static values are needed for $a$ and $b$. Appearances of the $f(x)$ shorthand can be replaced by the right-hand side of that equation. So, given a concrete linear function where $[ a \mapsto -1 , b \mapsto 5 ]$ we can calculate $f(4)$ as $-1 \cdot 4 + 5 = 1$.

\subsection{Area of Rectangle}
\label{sec:func:area}

We can also have function that are parameterized over two independent values. Such a function, for calculating the area of a rectangle, would look like this:
\begin{eqnarray*}
  area(\textcolor{purple}{w}, \textcolor{purple}{h}) = \textcolor{purple}{w} \cdot \textcolor{purple}{h}
\end{eqnarray*}

% description
The underlying principle is the same as it was for the linear function: By providing concrete values for each of the parameters, we have enough information to calculate the right-hand side of the equation and thus determine what the function evaluates to.

\section{\csharp}

% intro: recall that definition of f(x) from section \ref{sec:func:linear}? The csharp equivalent looks like this:
Recall that definition of $f(x) = -x+5$ from section \ref{sec:func:linear}? In order to create a \csharp\ equivalent, we need to map $x$ and $f(x)$ to specific types. Let's decide that both are \typename{double}. Then the \csharp\ implementation looks like this:

% example: linear function
\inputminted[linenos]{csharp}{../src/csharp/functions/f/Program.cs}

% followup and continuation
Lines 1-3 is where the action happens. This is where the function is defined. We will spend the remainder of this section going through it. For now, take a look at it and notice the similarities. The last three lines are simply there to \say{call} the \funcname{f} function with a number of different arguments and print out the result. This is to demonstrate how to use that definition. Let's do the same with the $area(w,h)$ function from section \ref{sec:func:area}:

% example area function
\inputminted[linenos]{csharp}{../src/csharp/functions/area/Program.cs}

% explain:
Again, the first three lines define a function. This time around, it has two parameters. But otherwise, it is largely the same. Lines 5-9 demonstrates how to use the function. Specifically, line 7 transfers execution to the function with two concrete values for width and height.

% syntax: in this section we will cover the definition and use of functions in csharp, specifically we will cover the elements from syntax \ref{syntax:fun}
In this section, we will cover the definition and use of functions in \csharp. Specifically, we will cover the elements from syntax \ref{syntax:fun}. It may look complicated, but you will quickly learn how to interpret it, and then it will become second nature.

\begin{syntaxfloat}
  \input{syntax/fun.tex}
  \caption{Functions.}
  \label{syntax:fun}
\end{syntaxfloat}

\subsection{Anatomy of a Function Definition}

% intro: recap of what a mathematical function is (mapping from input values to output value), figure \ref{fig:fun:anatomy} exemplifies the structure of the programming \say{equivalent}, equivalent is in quotation marks as \csharp function can do more than a pure mathematical functions. This is covered in section \ref{sec:func:csharp:sideeffects}. we will take some time to go through these in detail
At its core, a mathematical function maps a set of input values to an output value. Figure \ref{fig:fun:anatomy} exemplifies the structure of the \csharp\ \say{equivalent}. Here, quotation marks are used to signify that \csharp\ functions can do more than pure mathematical functions. What this entails is covered in section \ref{sec:func:csharp:sideeffects}. It will take some time to go through this figure in detail.

\begin{figure}[tbp]
  \begin{center}
\begin{minipage}{5.7cm}
\begin{minted}[breaklines,highlightlines={},escapeinside=||]{csharp}
|\tikzmark{returntypebefore}|int|\tikzmark{returntypeafter}| |\tikzmark{sigbefore}||\tikzmark{namebefore}|max|\tikzmark{nameafter}| (|\tikzmark{parbefore}|int num1, int num2|\tikzmark{parafter}|)|\tikzmark{sigafter}| {

|\tikzmark{bodybefore}|  int result;

if (num1 > num2) {
  result = num1;
} else {
  result = num2;
}

|\tikzmark{bodyafter}|  return result;|\tikzmark{returnstatement}|
}
\end{minted}
\end{minipage}
\end{center}
\begin{tikzpicture}[remember picture]
  % method name
  \draw[overlay, decorate,decoration={brace, amplitude=10pt, raise=11pt}] (pic cs:namebefore) to node[black,midway,above=21pt] {\rotatebox{90}{name}} (pic cs:nameafter);
  
  % parameter list
  \draw[overlay, decorate,decoration={brace, amplitude=10pt, raise=11pt}] (pic cs:parbefore) to node[black,midway,above=21pt] {parameter list} (pic cs:parafter);
  
  % signature
  \draw[overlay, decorate,decoration={brace, mirror, amplitude=10pt, raise=5pt, aspect=0.8}] (pic cs:sigbefore) to node[black,midway,below=15pt] {\hspace{28mm} signature} (pic cs:sigafter);
  
  % return type
  \draw[overlay, decorate,decoration={brace, amplitude=10pt, raise=11pt}] (pic cs:returntypebefore) to node[black,midway,above=21pt] {\rotatebox{90}{return}\\\rotatebox{90}{type}} (pic cs:returntypeafter);
  
  % prototype
  \draw[overlay, decorate,decoration={brace, mirror, amplitude=10pt, raise=5pt}] ([yshift=3mm] pic cs:returntypebefore) to node[black,midway,left=15pt] {\rotatebox{0}{function prototype}} ([yshift=-1mm] pic cs:returntypebefore);
  
  % return statement
  \draw[overlay, decorate,decoration={brace, amplitude=10pt, raise=5pt}] ([yshift=3.5mm] pic cs:returnstatement) to node[black,midway,right=15pt] {return statement} ([yshift=-2mm] pic cs:returnstatement);
  
  % body
  \draw[overlay, decorate,decoration={brace, mirror, amplitude=10pt, raise=5pt}] ([yshift=2mm] pic cs:bodybefore) to node[black,midway,left=15pt] {\rotatebox{90}{function body}} (pic cs:bodyafter);
\end{tikzpicture}

  \caption{Anatomy of a function definition.}
  \label{fig:fun:anatomy}
\end{figure}

% parameters vs arguments: a function is a parameterized piece of code, the concrete values that the parameters are bound to are called arguments, this constitutes the input interface of the function
A function is a parameterized piece of code. We can describe that function from the outside through its three interfaces, namely: The input interface, the output interface and the side-effect interface. On the input side, we have typed parameters. When invoked (or as we often say: called), the parameters available to the function body are bound to concrete values from the calling side. We can refer to these concrete values as \idx{arguments}{Argument!Function}. % TODO: revisit and spread out

% function interface: input, output, side-effect
We can describe the interface of a function from three sides:
\begin{enumerate}
  \descitem{Input Interface:} The input interface consists of the sequence of types of the function parameters.
  \descitem{Output Interface:} The output interface consists of the function return type.
  \descitem{Side-Effect Interface:} The side-effect interface is an umbrella term that covers all other interactions. This includes modifications to state external to the function and printouts to the screen. More on this in section \ref{sec:func:csharp:sideeffects}.
\end{enumerate}

% contract

\subsection{Parameter Passing}

% intro

% example
\inputminted[linenos]{csharp}{../src/csharp/functions/parampassing/Program.cs} % TODO: add parameter bindings

% explain

\subsubsection{Pass By Value vs Pass By Reference}

% intro

% example

% explain

\subsection{Returning Values}

% introduction: the output interface of the function, the explicit way of describing what (as in type) the function evaluates to
The output interface of a function is explicitly described through the return type. In the subset of \csharp\ that we cover in this book, only one value can be returned. That value can be complex though. Do note that this is an artificial restriction that we add to keep things (relatively) simple.

% naming functions after their type
If we want the function to evaluate to an \typename{int} then the return type should be \typename{int}. We can have any valid type as return type, even the complex types that we define ourselves. A call to a function can be placed anywhere in the code where a value of the functions return type can be placed. The return type is thus defining for where the function can be called. For that reason, we often refer to a function that return an \typename{int} as an \typename{int} function. And that convention is, of course, not limited to the \typename{int} type. Let's take a look at an example:

% example: int get_answer()
\inputminted[linenos]{csharp}{../src/csharp/functions/return/Program.cs}

% explain
In the first three declares one of the simplest possible functions that returns a value. Whenever it is called, it returns the value 42 as an \typename{int}. This stems from line two that contains a return statement. On line five, it is called. When executed, the program thus prints out \say{The answer is 42}. Note that line two is valid because it matches the last rule for \conceptname{stmt} in syntax \ref{syntax:fun}. Here, \valuename{42} plays the role of the expression, and it has a type that satisfies the return type of the function. This means that it either is an \typename{int} or it has a type which can be implicitly cast to an \typename{int}.

% the void type
Unlike mathematical functions, not all \csharp\ functions evaluate to a value. Perhaps they simply print out a datastructure in a readable way. Instead of having such functions return some value that doesn't carry any meaning (e.g., a zero) have have the return type follow from that value, \csharp\ has a \typename{void} type. If a function has a return type of \typename{void} then it doesn't evaluate to a value and is thus not allowed to return a value.

% example: void print_array()
\inputminted[linenos]{csharp}{../src/csharp/functions/void/Program.cs}

% explain
Again, we start out with a function definition followed by an application of it. This function takes an \typename{int[]} (i.e., and array of \typename{int} values) as parameter and doesn't return anything. That last part can be determined from it having a return type of \typename{void}. But, if it doesn't return anything, what is its function? On line three, it prints out all the elements of the array. It makes sure to add a space before each element. Notice the imbalance of spacing on lines two and four. This ensures that there will always be a single space padding between the first/last element and the square brackets. This little function could turn out quite convenient.

% return statement rules: every path through the function must end in a return statement, nothing must exist after a return statement (how would it be reached? indicator for a programming error), all return values must match the return type of the function
A special return statement is used to explicitly state which value to evaluate to. The valid uses of return statements are governed by the following rules:
\begin{enumerate}
  \item Every path through the function must end in a return statement (unless it is a \typename{void} function).
  \item A \typename{void} function has an implicit return statement at the end.
  \item No path through a function is allowed to have any statement after a return statement. How would such statement be reached? It is a clear indicator of a programming mistake.
  \item All return values must match the return type of the surrounding function. If it is a \typename{void} function, then the return statement does not take an expression.
\end{enumerate}

% when can a call be placed as an expression? when does it need to be as a statement?
Any function call can be used as a statement. When used as a statement, any return value will be discarded. For a function call to be used as an expression, it naturally must produce a value and that value must be of a type that is compatible with the use of that expression (i.e., its parent node in the parse tree). A \typename{void} function does not return a value and a call to it can thus never be used as an expression.

\subsection{Side-Effects}
\label{sec:func:csharp:sideeffects}

% the side-effect interface of the function: the primary interface consists of inputs and outputs, side-effects is an umbrella term that covers every other effect (that is, how it affects the part of the world that is external to it)
The primary interface of a function consists of its inputs and outputs. Functions can, however, also have \textsl{side-effects}. \idx{Side-effect}{Side-effect} is an umbrella term that covers every other effect. That is, how it affects the part of the world that is external to it. Examples of this include:
\begin{itemize}
  \item Modifying data that is not local to the function. This would be through a reference.
  \item Printing to the screen.
\end{itemize}

% example intro
Let's take a look at what this looks like in practice:

% example
\inputminted[]{csharp}{../src/csharp/functions/sideeffects/Program.cs}

% explain: let's first identify what is present in the program: a single global variable is defined then comes two functions and finally the "main" code that calls one function after the other, the first function (incr) increments value, it is allowed to do so as value is in global scope, updating a variable of a different scope than that of the function is a side-effect, the second function takes an integer as parameter, it then prints out a number of dashes that corresponds to that parameter and finally breaks the line, each of these printouts is a side-effect
Let's first identify what is present in the program: A single \idx{global variable}{Variable!Global} is defined, then comes two functions and finally the \say{main} code that calls one function after the other. The first function is called \funcname{incr} and it increments to value of the variable \varname{value}. It is allowed to do so, as the value is in \idx{global scope}{Scope!Global}. Updating a variable of a different scope than that of the function is a side-effect. The second function is called \funcname{print\_line}. It takes an integer as a parameter. It then prints out a number of dashes that corresponds o that parameter, and finally it \idx{breaks}{Linebreak} the line. Each of these printouts constitute a side-effect.

% function purity
Functions that neither have or rely on side-effects are said to be \idx{pure}{Purity!Function}. Such functions are \idx{stateless}{Stateless} and \idx{deterministic}{Deterministic} (in the sense that they will always produce the same output given the same input). This means that the function call can be substituted with the value it evaluates to. This quality is called \idx{referential transparency}{Transparency!Referential}.

\subsection{Caller and Callee Roles}

We are always in some function, even in our hello world example from section \ref{first:projectcreation}. The entirety of this example is a single statement. The compiler will essentially wrap that code in a dummy function. So, whenever a function is called, there are two roles: The caller and the callee. If the function \funcname{a} calls the function \funcname{b}, then we refer to \funcname{a} as the \idx{caller}{Caller} and \funcname{b} as the \idx{callee}{Callee}.

% example
\inputminted[]{csharp}{../src/csharp/functions/roles/Program.cs}

% followup: note that the role is associated with a specific function call, imagine that we also had a \funcname{c} function so that a called b which called c which simply returned its parameter, then b would be the callee of the first call from a and the caller of the second call to c
Note that the role is associated with a specific function call. Imagine that we also had a function \funcname{c} such that \funcname{a} called \funcname{b} which called \funcname{c} which in turn just returned the value of its parameter. Then \funcname{b} would be the callee of the first call from \funcname{a} as well as the caller of the second call to \funcname{c}.

\subsection{Dispatch}

% problem: in order to call a function a concrete implementation must first be identified, a dispatch mechanism is then responsible for handing over the thread of execution to the callee
In order to call a function, a concrete implementation must first be identified. This is done for us by the compiler. A \idx{dispatch mechanism}{Dispatch mechanism} is responsible for handing over the \idx{thread of execution}{Thread of execution} to the called function.

% solution: in \csharp this lookup is done based on the combination of the function name and the ordered sequence of parameter types, as a consequence we are allowed to have two function definitions under the same name as long as their parameter type lists differ, that is their signatures are different, if we have two identical signatures then \csharp wouldn't be able to differentiate between them and thus be unable to uniquely dispatch in a consistent manner, because of this such code will result in a compile-time error
In \csharp, this lookup is done purely based on the function name. As a consequence of this, we are not allowed to have two functions of the same name. \csharp\ wouldn't be able to differentiate between them and thus be unable to uniquely dispatch in a consistent manner. Because of this, such code will result in a \idx{compile-time error}{Error!Compile-time}.

% example

% explain

\subsection{The Stack}

% introduction: as the body of a function is an ordinary statement functions can call functions, this leads to one challenge in particular: as the code of a function is compiled to machine code variables are mapped to registers, what happens if a function calls itself? Are those variables then shared across the two invocations? while that could be useful it would certainly be confusing, in fact most programming language designers (those of \csharp included) have decided that it too confusing, their solution is to include a program-wide datastructure called *the stack* that can store the registers of the calling function so that the called function is free to do with them as it pleases, as a side-effect this greatly simplifies the compilation process as compilation can be done on a per-function basis
As the body of a function is an ordinary statement, functions can call functions. This leads to one challenge in particular: As the code of a function is compiled to \idx{machine code}{Machine code}, \idx{variables}{Variable} are mapped to \idx{registers}{Register}. That happens if a function calls itself? Are those variables then shares across the two \idx{function invocations}{Function!Invocation}? While that could be useful in some situations, it would certainly be confusing in most. In fact, most programming language designers (those of \csharp\ included) have decided that it is too confusing. Their solution is to include a program-wide datastructure called \say{\idx{the stack}{Stack!The}} that can store the registers of the calling function so that the called function is free to do with them as it pleases. The reality is a bit more complicated, but at this point it doesn't matter. As a side-effect, this greatly simplifies the \idx{compilation}{Compilation} process as compilation can be done on a per-function basis.

% what is a stack: think of a stack as a stack of papers, a stack is a form of data structure that has two operations: you can push a value to the top of the stack, and you can remove a value from the top of the stack. That last operation is called a "pop" operation on the stack, in reality this means that a stack operates as a last-in-first-out (or LIFO) mechanism, stacks are generic datastructures that are frequently used in programming, when we refer to "the stack" then it is the stack that keeps track of the variables that are local to calling function invocations
Think of a stack as a stack of papers. It is a form of data structure that has two operations: You can add a value to the top of it (through a \idx{push operation}{Push operation}) and you can remove a value from the top of it (through a \idx{pop operation}{Pop operation}). In reality, this means that the stack operates as a last-in-first-out (or \idx{LIFO}{LIFO}) mechanism. Stacks are generic datastructures that are frequently used in programming. When we refer to \say{\idx{the stack}{Stack!The}} then it is the stack that keeps track of the variables that are local to calling function invocations that we mean.

% figure: stack

% stack frames: instead of placing individual values on the stack we place stack frames, these are structures of data that match what needs to be saved, everything that needs to be reestablished on return from the function call can be "removed" in a single operation
Instead of placing individual values on the stack, we place \idx{stack frames}{Stack frame}. These are structures of data that match what needs to be saved. Everything that needs to be reestablished on return from the function call can then be popped in a single operation.

% TODO: which parts of this are generic enough to warrant being placed before the csharp section?

% example

% explain

\subsection{Recursion}

% fibonacci: we will start this section with a minor deviation, in mathematics the fibonacci sequence is a sequence of integer numbers where each element is the sum of the two previous ones, the sequence was originally described by indian mathematicians (e.g., Acharya Pingala [पिङ्गल] around 200 BC) but today we know it through the name of fibonacci (Leonardo of pisa) who rediscovered the sequence in 1202, this particular sequence appear unexpectedly often in mathematics and nature, and is related to the golden ratio, it has even found its way into music, formally we can define the sequence as
We will start this section with a minor digression. In mathematics, the \defi{Fibonacci sequence}{Fibonacci sequence} was originally described by Indian mathematicians (e.g., by \idx{Acharya Pingala}{Pingala, Acharya पिङ्गल} -- पिङ्गल in sanskrit-- around 200 BC\cite{towards_a_global_science}) but today we know it through the name of Fibonacci (aka \idx{Leonardo of Pisa}{Leonardo of Pisa}) who rediscovered it in 1202. This particular sequence appears unexpectedly often in mathematics and nature, and is related to the \idx{golden ratio}{Golden ratio}. It has even found its way into music\cite{tool_lateralus}. Formally, we can define the $n$'th element of the sequence as:

% formula
\begin{eqnarray}
  \nonumber
  fib(n) =
    \begin{cases}
      0, & \text{\color{violet} for $n=0$} \\
      n, & \text{\color{violet} for $n=1$} \\
      fib(n-1)+fib(n-2), & \text{\color{violet} for $n>1$} \\
    \end{cases}
\end{eqnarray}

% interface: What does the equivalent look like in csharp code? let's start by defining the interface of the function that we are about to write, it is obvious that it takes one parameter (the $n$), the n refers to an index for the sequence, this indicates that it is an unsigned integer, that unsigned integer constitutes the input interface of the function, the output is 0 (some integer) n (same as input) or the result of adding two of these, adding a unsigned value and a zero will never yield a negative value and so the output interface must be an unsigned integer, there does not appear to be a need for it to print out anything and it doesn't operate on state besides its input, that means that the side-effect interface is empty, so, we are looking for a function with the prototype: uint fib (uint n)
What does the \csharp\ equivalent look like in code? Let us start by defining the interface of the function that we are about to write. It should be obvious that it takes one parameter (the $n$). That $n$ refers to the index of a sequence, and that indicates that it is an \idx{unsigned}{Integer!Unsigned} integer. This unsigned integer constitutes the input interface of our function. The output is zero (some integer), $n$ (unsigned integer as input) or the result of adding two of these. Adding an unsigned value and a zero (or another unsigned integer for that matter) will never yield a negative value, and so the output interface must be an unsigned integer. There does not appear to be a need for it to print out anything, and it doesn't operate on state besides its input. That means that the side-effect interface is empty. So, we are looking for a function with the \idx{prototype}{Prototype!Of function}:
\begin{center}
  \texttt{\typename{uint} \funcname{fib} (\typename{uint} \varname{n})}.
\end{center}

% intro: looking at this definition there is no obvious way of implementing the function using branches loops and our basic datastructures, but once you \textsl{get} the idea it is trivial to implement using function calls, that is function calls to itself, one way of doing so looks like this:
Now that we understand the interface of the function, let's take a look at its body. Looking at the above definition, there is no obvious way of implementing the \funcname{fib} function using branches, loops and our basic datastructures. But once you \textsl{get} the idea it is trivial to do so using function calls. That is, function calls to itself. One way of doing so looks like this:

% example
\inputminted[]{csharp}{../src/csharp/functions/fib/Program.cs}

% explain: there is a one-to-one mapping between the mathematical formula and this implementation, a switch construct is used to differentiate between the three cases of the formula, the default matches >1 as 0 and 1 have been covered by the previous cases and n is of an unsigned integer type: it can never represent a negative number
There is a one-to-one mapping between the mathematical formula and this implementation. A \keywordname{switch} construct is used to differentiate between the three different cases of the formula. The \textsl{default} matches $n>1$ as $n=0$ and $n=1$ have been covered by previous cases and \varname{n} is of an \idx{unsigned}{Integer!Unsigned} integer type: It can never represent a negative number.

% pattern: but how does it work? on the surface we have taken one problem (fib(n)) and converted it to two problems (fib(n-1) and fib(n-2)), while this may not seem like progress each of these subproblems are easier to solve than the original, for instance for n=2 this will convert to fib(1) and fib(0), each of these are trivial and fib(2) is thus a bit harder than trivial, and fib(3) is a bit harder yet again, there is an underlying pattern here: we have one or more trivial base cases and a recursion step that solution-wise brings us closer to those base cases. When a function calls itself in order to solve a subproblem we call the function recursive, recursion is a form of repetition and can -- in some cases -- be a valid or even superior alternative to loops

\subsubsection{Computation}

\subsubsection{Stack}

\subsection{Guards}

% intro: consider the following problem: you are hosting a children's birthday party and you have baked a nice round cake, you want all participants to receive a slice of equal size, how many degrees should each slice be?
Consider the following problem: You are hosting a children's birthday party, and you have baked a nice round cake. You want all the participants to receive a slice of equal size. How many degrees should each slice be?

% transfer to code: what would this look like in code? we have to deliver an angle and that angle depends on a number of participants, this only makes sense for a positive number of participants, so we want to make sure to treat the cases of negative and zero counts in a special way, all things considered, this is pretty simple, we can try out a few different inputs:
What would this look like in code? We have to deliver an angle, and that angle depends on a number of participants. Note that this only makes sense for a positive number of participants, so we want to make sure to treat cases of negative and zero counts in a special way. We don't have a type that only supports positive integers. The closest we can get is to use a non-negative integer type, and that would still leave us with a \idx{special case}{Special case}. But all things considered, this is pretty simple. We can try out the functionality with a few different inputs:

% example: no function example
\inputminted[]{csharp}{../src/csharp/functions/no_guard_no_function/Program.cs} % TODO: align variable names to story

When executed, this program prints:
\begin{minted}{shell-session}
$ dotnet run
Slices of 360 degrees will feed 1 mouths
Slices of 180 degrees will feed 2 mouths
Slices of 120 degrees will feed 3 mouths
Slices of 90 degrees will feed 4 mouths
\end{minted}

% adding the function: the code works but the logic of calculating the angles is mixed in with the logic for cycling through the different inputs, understaning one of these two functionalities requires an overview of both functionalities, in a program of this size that is not a problem but that changes once we get to play with larger codebases, to deal with this we need to find highly coherent parts of the code that can be parameterized as functions and thus abstracted away, in this case the natural abstraction is a function that given a number of participants calculates an angle, then that function can be called for each of the numbers of participants, this allows us to reason about each of the two parts without understanding the intricacies of the other, that looks like this:
The code works, but the logic of calculating the angles is mixed in witht the logic for cycling throught the different inputs. It is hard to tell where one starts and the other ends. Understanding one of these two functionalities requires an overview of both functionalities. In a program of this size, that is not a problem, but that changes once we get to play with larger codebases. To deal with this, we need to find highly coherent parts of the code that can be \idx{parameterized}{Paramterization} as functions and thus be \idx{abstracted}{Abstraction} away as a function call. In this case, the natural abstraction is a function that, given a number of participants, calculates an angle. Then that function can be called for each of the numbers of participants. This allows us to reason about each of the two parts without understanding the intricacies of the other. It looks like this:

% example: function without guard (something like splitting a cake into n pieces)
\inputminted[linenos=true]{csharp}{../src/csharp/functions/no_guard/Program.cs}

% observations: let's look at that slice_angle function and in particular line 5, in this particular example that line is trivial: it returns the value -1, but imagine that we wanted to have some other code in his block, we know that mouth_count is positive but why do we know that? we know that because we are in the else branch of an if-else construct whose condition tells us so, that condition is expressed higher up in our code, only by two lines but still, this distance represents a context window that we as programmers have to spend mental resources keeping fresh in mind, the larger and more complex the function the larger the burden, let's see if we can simplify it:
Let's take a look at that \funcname{slice\_angle} function, and in particular line 5. At first glance, this is trivial: It returns the value \valuename{-1}. But imagine that we wanted to have some other code in this block. We know that \varname{mouth\_count} is positive, but why do we know that? We know that because we are in the \keywordname{else} branch of an \keywordname{if-else} construct whose conditions tells us so. That condition is expressed higher up in the code. Only by two lines, but still \ldots\ This distance represents a \idx{context window}{Context window!Mental} that we as programmers have to spend \idx{mental resources}{Resource!Mental} on in order to keep it fresh in mind. The larger and more complex the function, the larger the burden. Let's see if we can simplify it:

% example: same code with a guard
\inputminted[lastline=7,linenos=true]{csharp}{../src/csharp/functions/guard/Program.cs}

% observations and expansion: first of all this implementation does the same as the previous, in fact it will likely be compiled to the exact same code, in its previous incarnation this function split out the two cases using an if-else construct, all code were either in the true block or in the false block, in the new version there is no false block, or rather because the true block returns the remainder of the function in essense becomes the false block, lines 2-4 are there to filter out a special case, after line 4 we no longer care about what happened up there because we know that our input is sane and from here on it will be smooth sailing
First of all, this implementation is \idx{functionally equivalent}{Equivalency!Functional} to the previous. In fact, it will likely be \idx{compiled}{Compilation} to the exact same code. In its previous incarnation, this function split out the two cases using an \keywordname{if-else} construct. There was no code outside of the \keywordname{true} block or the \keywordname{false} block. In the new version, there is no \keywordname{false} block. Or rather, because the \keywordname{true} block returns, the remainder of the function in essense becomes the \keywordname{false} block. Lines 2-4 are there to filter out a \idx{special case}{Special case}. After line 4, we no longer case about what happened up there. We don't need to. All we need to know is that our input is sane, and from there on it will be smooth sailing. This is called a \idx{guard}{Guard}.

% outro: often you will see a small series of such guards, each rejecting a special case, and then comes the actual implementation of the function, it is good style to lead each guard with a comment explaining (i) that it is a guard, and (ii) which guarantees it provides
Often, you will see a small series of such guards -- each rejecting a spacial case -- and then comes the actual implementation of the function. It is good style to lead each guard with a comment explaining (i) that it is a guard, and (ii) which \idx{guarantees}{Guarantee} it provides.

\subsection{Nesting}

% intro: functions can be nested, what that means, why would you do this? structure (if the function is only meant to be called from within this particular other function), scope nesting (if the function is simpler to write if it gets access to the function-local scope of the calling function), let's take a look at an example:
Functions can be \defi{nested}{Nesting!Of functions}. This means that a function can be defined inside the body of another function definition. But why would you do this? The \idx{scope}{Scope} of the inner function will will be limited to the body of the outer function. This has two important (potential) benefits: First, If a function is only ever meant to be called from within a particular other function then by placing it within the body of this other function we remove the possibility of it being called from anywhere else. Second, if this function interacts with lots of variables local to another function then by placing it insite that functions body, it will gain access to the \idx{function-local}{Scope!Function-local} scope and we won't need to send references to and from the function call. This could make the interface a lot simpler. There can be lots of downsides to nesting function definitions, and often it simply isn't an option. The most commonly cited downside is perhaps that it is not a common thing to do, and thus it will have a negative impact on code \idx{readability}{Readability} for a large fraction of developers.

% example
\inputminted[linenos=true]{csharp}{../src/csharp/functions/nesting/Program.cs}

% explain:
Had it not been for lines five through nine you would likely not had bat an eye at this. To be clear, no sane programmer would write the code like this; it is unnescesarily \idx{complex}{Complexity}. But it does allow us to discuss nested functions. In this case, the \funcname{operation} function is defined within the scope of the body of the \funcname{fun}. For every \idx{invocation}{Invocation!of function} of the \funcname{fun} function, there are two local variables, namely \varname{a} and \varname{b}. When the \funcname{operation} function is invoked, it is always within this scope, and thus it is capable of naming these variables (aka it has access to them). When this program is executed, it prints out \say{\texttt{3}}.

\section{Python}

% intro: python is a scripting language, these are generally weakly typed and so is python, this means what (type boundaries are not strictly enforced, variable names are not restricted to specific typed on declaration and often not declared at all, over time the type of the variable may change or you could say that it can name values of different type, implicit type conversion), accordingly function declarations do not declare parameter or return types, in python functions are declared using the def keyword:
\idx{Python}{Python programming language} is a \idx{scripting language}{Language!Scripting}. These are generally weakly typed, and that is also the case for Python. \defi{Weakly typed}{Typing!Weak} languages are characterized by having no strict enforcement of \idx{type boundaries}{Type!Boundaries}. This means that variable names are not restrictted to a specific type on declaration (if they are ever declared) and over time the type of the variable may change (or you could say that it can name values of different type),. These language also tend to be fond of \idx{implicit type conversion}{Type conversion!Implicit}. Function declarations in Python thus don't declare parameter or return type. Let's take a look at such a declaration:
% TODO: move part of this to the primitives.tex file?

% example: add function
\inputminted[firstline=3]{python}{../src/python/functions/add/add.py}

% there is no separate compilation phase, function lookup is dynamic, we can get a runtime error from a missing function
As is common with scripting languages, there is no separate compilation phase and \idx{function lookup}{Function lookup} is \idx{dynamic}{Function lookup!Dynamic} (i.e., it happens at \idx{runtime}{Runtime}). Because of this, if we are missing a function declaration for a specific call, that will be discovered when we try to evaluate that call. This results in a \idx{runtime error}{Error!At runtime}, and it can take some time and specific inputs for it to appear.

\section{Elixir}

% anonymous functions, this is a concept that is rooted in the functional paradigm, we will see real functions later on but in this section we focus on anonymous functions, but how do you call them if they are nameless? well while they are not assigned a static name they are values and can thus be assigned to variables, those variables have names and through those names we can call them:
\idx{Elixir}{Elixir programming language} is an example of a \idx{functional language}{Functional programming language}. One of the key element of the functional \idx{paradigm}{Paradigm!Functional} is that we have access to \idx{anonymous functions}{Function!Anonymous}. Such functions are anonymous is the sense that the are not assigned a static \idx{name}{Name} thought which they can be called. But how can we call them if they are nameless? Well, in functional languages, functions are values and can thus be assigned to variables. Those variable have names, and throught those names we can call them:

% example: anonymous function definition and call
\inputminted[firstline=3]{elixir}{../src/elixir/functions/anon/anon.exs} % TODO: Convert to iex

% explain: the first line declares an anonymous function and assigns it to a variable called incr, on the right-hand side of the assignment we find the declaration, the parameter name is found between the fn keyword and the arrow, and the function body is between the arrow and the end keyword, in elixir the function evaluates to whatever the last statement of a function body evaluates to, this works because everything is an expression in elixir, on the second line this function is called, note the period operator, this tells the compiler that we want to treat the variable as a function and call it. This will evaluate to 12+1=13.
The first line declares an anonymous function and assigns it to a variable called \varname{incr}. On the \idx{right-hand side}{Right-hand side} of the assignment, we find the declaration. The parameter name is found between the \keywordname{fn} keyword and the \keywordname{->} keyword, and the function body is betwen the \keywordname{->} keyword and the \keywordname{end} keyword. In Elixir, functions evaluate to whatever the last statement of their function body evaluates to. This works because everything is an \idx{expression}{Expression} in Elixir. On the second line, this function is called. Note the \keywordname{.} operator. This tells the compiler that we want to treat the variable as a function and call it. This will evaluate to $12+1=13$.

\subsection{Pipelines}

% pipelining: we often see the result of a function call be used only as an argument for another function call, this pattern is even more common in functional programming, elixir provides syntactic sugar for this pattern in the form of the \keywordname{|>} operator, it injects the result of the function call on its left-hand side as the first argument of a function call on its right-hand side. Because of this, the latter function call will appear to be missing an argument in the code. The elixir standard library is full of functions that are designed to exploit this. But let's take a look at a concrete example:
We often see the result of a function call be used only as an argument for another function call in the next statement. This \idx{pattern}{Pattern} is even more common in \idx{functional programing}{Programming!Functional} than it is in \idx{imperative programming}{Programming!Imperative}. Elixir provides \idx{syntactic sugar}{Syntactic sugar} for this pattern in the form of the \keywordname{|>} operator (aka the pipe operator). It injects the result of the function call in its \idx{left-hand side}{Left-hand side} as the first argument of the function call on its right-hand side. Because of this, the latter function call will appear to be missing an argument in the code. But that is not the case. The Elixir \idx{standard library}{Standard library} if full of functions that are designed to exploit this pattern. But let's take a look at a concrete example:

% example
\inputminted[firstline=4]{elixir}{../src/elixir/functions/pipelining/pipelining.exs} % TODO: Convert to iex

% explanation: double application of the pipe operator, first expression (11) given as input to incr, evaluates to 12, given as input to incr, evaluates to 13
We have two applications of the pipe operator. This creates a \say{pipeline}. Fist, the expression \valuename{11} is given as argument to the first call to \funcname{incr}. This evaluates to $11+1=12$. That 12 is then given as input to the second call to \funcname{incr} which evaluates to $12+1=13$.

\subsection{Higher-Order Functions}

% higher-order functions: there are a number of side-effects from functions being values, one such side-effect is that a function can take a function as parameter and have a function as return value, It can also construct a new function dynamically and return it, functions that have other functions as part as their parameters or return values are called higher-order functions, a few common patters exist, lets start with the map function, explain what it does, let's see this in action:
There are a number of side-effects from functions being \idx{values}{Value}. One such effect is that a function can take functions as \idx{parameters}{Function!as parameter} and have a function as \idx{return value}{Function!as return value}. It can also construct a new function dynamically and return it. Functions that have other functions as part of their parameters or return values are called \defi{higher-order functions}{Function!Higher-order}. We often say that these functions are \idx{composable}{Composability!of functions}. This property has resulted in a number of \idx{patterns}{Pattern} to be defined. Let's start with the \defi{map}{Pattern!Map} pattern. In this pattern a function is applied to each element of a list to construct a new list. In action, it looks like this:

% example: the map function
\inputminted[firstline=5]{elixir}{../src/elixir/functions/map/map.exs} % TODO: How to get rid of the awkward printout

% explain: again a pipeline with two steps, input is a list with all elements from 1 to 10 (both ends included), then we use the map function to create a new list of the same size where the incr funciton is applied to each value, finally we inspect the result by printing it out, due to a leftover from erlang printing out a list of small integers is a bit intricate, just accept that the last line prints out the result
Again, we have a pipeline in three phases. Input is a list with all integers from one to ten (both ends included). In the next phase, the \funcname{Enum.map} function is used to create a new list of the same size where the \funcname{incr} function is applied to each of the corresponding values in the input list. Finally, we inspect the result by printing it out. Due to a leftover oddity from \idx{Erlang}{Erlang programming language}, printing out a list of small integers is overly complicated. Just accept that this is what the last line does.

\subsection{Larger Example}

% filter and reduce: the map function should be seen as a pattern where something is done to every element of a list, there are other such patterns, the filter pattern constructs a new list that retains all elements of an input list for which some function evaluates to true, explain reduce, the following is a more complex example that illustrate how expressive these can be:
The map function should be seen as a pattern where something is done to every element of a list. There are other such patterns. The \defi{filter}{Pattern!Filter} pattern constructs a new list that retains all elements of an input list for which some function evaluates to true. The \defi{reduce}{Pattern!Reduce} pattern repeatedly applies a function to each element and an \idx{accumulator value}{Accumulator value} (that it then updates). This creates one value from the contents of the entire input list. A simple use-case is to find the largest number in a list. The following code is a more complex example that illustrates how \idx{expressive}{Expressiveness} higher-order functions like these can be:

% example: larger higher-order example with map, filter and reduce
\inputminted[firstline=4]{elixir}{../src/elixir/functions/large/large.exs} % TODO: How to get rid of the awkward printout

% explain: so lets take a look at this mouthful, on the first line we define an function called positive, it can tell us whether an input is positive, then we run a pipeline, it starts out with a list of integers, some negative some zero and some positive, then it maps the incr function to them thereby incrementing each number by the value 1, next a filtering is done to keep the positive values, finally (if we disregard the printout) we reduce the remaining numbers, this is where it gets tricky to follow
So, let's take a look at this mouthful. One the first line, we define a function called \funcname{positive}. It can tell us whether an input is positive. Then we run a pipeline. It starts out with a list of integers. Som are negative, some are zere and some a positive. Then, it maps the \funcname{incr} function to them, thereby incrementing each number by the value 1. Next, a filtering is done to keep the positive values. Finally -- if we disregard the printout -- we reduce the remaining numbers to find the largest value. This is where it gets tricky to follow.

% what happens in the reduce: just like the map function applies a function to each element of its input so does the reduce function, the reduce function however takes two parameters, the second parameter is a so-called accumulator value, it carries a result from the last application of the funtion, an initial value (0 in this case) is applied for the first call to the "reducer" function, to follow the example we first find the max of the first element of the input list and zero, then we find the max of the result of that and the second element of the list, and so we continue until we have reached the end of the list and thus found the largest element, do note that -- just as for map and filter -- we can provide an arbitrary function (of the right interface) as parameter
Just like the \funcname{map} function applies a function to each element of its input list, so does the \funcname{reduce} function. That function, however, takes two parameters. The second parameter is a so-called \say{accumulator} value. It carries a result from the last application of the function. An initial value (\valuename{0} in the case of the example) is applied for the first call to the \say{reducer} function. To follow the example, we first find the \idx{max}{Maximum} of the first element of the input list and zero. Then, we find the max of the result of that and the second element of the list. And so, we continue until we have reached the end of the list, and thus found the largest element. Do note, that -- just as for \funcname{map} and \funcname{filter} -- we can provide an arbitrary function (of the right interface) as the parameter.

\subsection{Function Lookup}

% function lookup in the BEAM: elixir inherits the BEAM runtime from erlang, in it function dispatch is performed at runtime using function name and number of parameters, this mechanism can be used to redefine the implementation at runtime for hot-code updates, this trades function call speed for maintainability
\idx{Elixir}{Elixir programming language} inherits the \idx{BEAM}{BEAM runtime} runtime from \idx{Erlang}{Erlang programming language}. In it, \idx{function dispatch}{Function dispatch} is performed by lookup against a table at runtime using function name and parameter count. This indirection can be used to redefine the implementation of functions at runtime, and that ability is used to issue \defi{hot updates}{Hot update}. These are code updates to a running system. So, function call speed is traded for \idx{maintainability}{Maintainability}. This is especially important for long-running \idx{deployments}{Deployment}.

\exercises{functions}{Functions}
